\documentclass[10pt,a4paper]{amsart}
\usepackage{amsmath,amssymb,amsthm,graphicx}
\usepackage{verbatim}
\newenvironment{reference}{\comment}{\endcomment}
%\newenvironment{reference}{}{}
\newenvironment{slogan}{\comment}{\endcomment}
\newenvironment{history}{\comment}{\endcomment}

% For commutative diagrams we use Xy-pic
\usepackage[all]{xy}

% We use 2cell for 2-commutative diagrams.
\xyoption{2cell}
% Theorem environments.
%
\theoremstyle{plain}
\newtheorem{theorem}[subsection]{Theorem}
\newtheorem{proposition}[subsection]{Proposition}
\newtheorem{lemma}[subsection]{Lemma}

\theoremstyle{definition}
\newtheorem{definition}[subsection]{Definition}
\newtheorem{example}[subsection]{Example}
\newtheorem{exercise}[subsection]{Exercise}
\newtheorem{situation}[subsection]{Situation}

\theoremstyle{remark}
\newtheorem{remark}[subsection]{Remark}
\newtheorem{remarks}[subsection]{Remarks}

\numberwithin{equation}{subsection}

% Macros
%
\def\lim{\mathop{\mathrm{lim}}\nolimits}
\def\colim{\mathop{\mathrm{colim}}\nolimits}
\def\Spec{\mathop{\mathrm{Spec}}}
\def\Hom{\mathop{\mathrm{Hom}}\nolimits}
\def\Ext{\mathop{\mathrm{Ext}}\nolimits}
\def\SheafHom{\mathop{\mathcal{H}\!\mathit{om}}\nolimits}
\def\SheafExt{\mathop{\mathcal{E}\!\mathit{xt}}\nolimits}
\def\Sch{\mathit{Sch}}
\def\Mor{\mathop{\mathrm{Mor}}\nolimits}
\def\Ob{\mathop{\mathrm{Ob}}\nolimits}
\def\Sh{\mathop{\mathit{Sh}}\nolimits}
\def\NL{\mathop{N\!L}\nolimits}
\def\CH{\mathop{\mathrm{CH}}\nolimits}
\def\proetale{{pro\text{-}\acute{e}tale}}
\def\etale{{\acute{e}tale}}
\def\QCoh{\mathit{QCoh}}
\def\Ker{\mathop{\mathrm{Ker}}}
\def\Im{\mathop{\mathrm{Im}}}
\def\Coker{\mathop{\mathrm{Coker}}}
\def\Coim{\mathop{\mathrm{Coim}}}

% Boxtimes
%
\DeclareMathSymbol{\boxtimes}{\mathbin}{AMSa}{"02}

%
% Macros for moduli stacks/spaces
%
\def\QCohstack{\mathcal{QC}\!\mathit{oh}}
\def\Cohstack{\mathcal{C}\!\mathit{oh}}
\def\Spacesstack{\mathcal{S}\!\mathit{paces}}
\def\Quotfunctor{\mathrm{Quot}}
\def\Hilbfunctor{\mathrm{Hilb}}
\def\Curvesstack{\mathcal{C}\!\mathit{urves}}
\def\Polarizedstack{\mathcal{P}\!\mathit{olarized}}
\def\Complexesstack{\mathcal{C}\!\mathit{omplexes}}
% \Pic is the operator that assigns to X its picard group, usage \Pic(X)
% \Picardstack_{X/B} denotes the Picard stack of X over B
% \Picardfunctor_{X/B} denotes the Picard functor of X over B
\def\Pic{\mathop{\mathrm{Pic}}\nolimits}
\def\Picardstack{\mathcal{P}\!\mathit{ic}}
\def\Picardfunctor{\mathrm{Pic}}
\def\Deformationcategory{\mathcal{D}\!\mathit{ef}}
\title{A cartesian square induces an isomorphism on the kernels of its horizontal maps}
\author{The Stacks Project Authors}
\date{stacks-project:homology.tex:local-snapshot-20261012; GNU FDL 1.2 or later}
\begin{document}
\maketitle
\noindent\textit{Editorial note.} All delivered mathematical text is the author's own, copied from the pinned Stacks Project chapter file homology.tex; no statement or proof step was rewritten, reordered or removed.\par\medskip
\noindent\textit{Editorial note.} The proof is the author's complete argument for the cartesian statement, with both inverse relations computed and both cancellations justified by naming the monomorphism used; the closing clause 'Now, (2) follows by duality.' is the author's own routine duality remark on the part written out in full inside the counted span, and it is preserved verbatim.\par\medskip
\noindent\textit{Editorial note.} Every labelled reference occurring inside the counted spans is carried in assembly.external\_references as a non-load-bearing navigation pointer (same-chapter ones at their definition site, cross-chapter ones at the invocation inside the counted span), so no load-bearing mathematics is replaced by a stub. No bibliography citation occurs inside the counted spans.\par\medskip
\noindent\textit{Editorial note.} The same-name check for this card scanned every staged Stacks chapter of this repository -- all sources/stacks-*/pinned-original-source.tex files, the card's own chapter excluded -- and the labels found are recorded in dedup.same\_name\_cross\_chapter on the card.\par\medskip
\noindent\textit{Editorial note.} The preamble is the chapter's own native preamble excerpt already used for this source: the comment-type reference, slogan and history environments (native preamble lines 8-13), the Xy-pic setup (native preamble lines 14-18) and the native theorem declarations with the mathematical macros (native preamble lines 154-217, the actual last line of the file).\par\medskip
\noindent\textit{Allowed standard prerequisites (editorial).}\begin{itemize}
\item Cartesian squares as fibre products. Same chapter, the description of a cartesian square by the universal property of the fibre product. Supplies the morphism from the kernel of the lower horizontal map into the upper left object used at the start of the proof.
\item Kernel maps are monomorphisms. Same chapter, Lemma lemma-kernel-mono. Supplies the cancellation of the two kernel injections that turns the computed identities into the two required inverse relations.
\end{itemize}\medskip
\begin{lemma}
\label{lemma-cartesian-kernel}
Let $\mathcal{A}$ be an abelian category. Let
$$
\xymatrix{
w\ar[r]^f\ar[d]_g
& y\ar[d]^h\\
x\ar[r]^k
& z
}
$$
be a commutative diagram.
\begin{enumerate}
\item If the diagram is cartesian, then the morphism 
$\Ker(f)\to\Ker(k)$ induced by $g$ is an isomorphism.
\item If the diagram is cocartesian, then the morphism 
$\Coker(f)\to\Coker(k)$ induced by $h$ is an isomorphism.
\end{enumerate}
\end{lemma}

\begin{proof}
Suppose the diagram is cartesian. Let 
$e:\Ker(f)\to\Ker(k)$ be induced by $g$. Let 
$i:\Ker(f)\to w$ and $j:\Ker(k)\to x$ be the canonical 
injections. There exists $t:\Ker(k)\to w$ with $f\circ t=0$ 
and $g\circ t=j$. Hence, there exists $u:\Ker(k)\to\Ker(f)$ 
with $i\circ u=t$. It follows 
$g\circ i\circ u\circ e=g\circ t\circ e=j\circ e=g\circ i$ and 
$f\circ i\circ u\circ e=0=f\circ i$, hence $i\circ u\circ e=i$. Since 
$i$ is a monomorphism this implies $u\circ e=\text{id}_{\Ker(f)}$.
Furthermore, we have $j\circ e\circ u=g\circ i\circ u=g\circ t=j$. 
Since $j$ is a monomorphism this implies $e\circ u=\text{id}_{\Ker(k)}$.
This proves (1). Now, (2) follows by duality.
\end{proof}

\end{document}
